% 张三慧大学物理学（第4版）V2第23章 光的衍射
% 主教材PDF物理页287--289；辅导S PDF物理页321--325。
% 原图23.23及23.31--23.34需独立矢量重绘与复核。
\section*{第23章\quad 光的衍射}
\subsection*{章末习题23.1--23.19}
\paragraph{习题 23.1}
单缝宽 \(a=0.10\,\mathrm{mm}\)，其后有焦距 \(f=50\,\mathrm{cm}\) 会聚透镜，\(\lambda=546.1\,\mathrm{nm}\) 的平行光正入射。求焦平面中央明纹全宽。
\paragraph{习题 23.2}
波长 \(632.8\,\mathrm{nm}\) 的激光正入射单缝，夫琅禾费第一极小与法线夹角为 \(5^\circ\)。求缝宽。
\paragraph{习题 23.3}
未知单色光正入射单缝的第3级明纹，恰与 \(\lambda=600\,\mathrm{nm}\) 的光照射该缝时的第2级明纹重合。求未知波长，说明次级明纹位置的近似条件。
\paragraph{习题 23.4}
海面波长 \(20\,\mathrm m\)，垂直传入宽 \(50\,\mathrm m\) 的港口。求港内衍射中央波束的全角宽度（按单缝模型）。
\paragraph{习题 23.5}
用肉眼观察星体，瞳孔最大直径 \(7.0\,\mathrm{mm}\)，光波长 \(550\,\mathrm{nm}\)。(1) 求艾里斑的全角宽；(2) 瞳孔距视网膜 \(23\,\mathrm{mm}\)，求斑直径；(3) 中央小凹直径 \(0.25\,\mathrm{mm}\)，感光细胞面密度 \(1.5\times10^5\,\mathrm{mm^{-2}}\)，估算星像覆盖的细胞数。
\paragraph{习题 23.6}
在离地 \(3.5\times10^4\,\mathrm{km}\) 的卫星上发射波长 \(10\,\mathrm{cm}\) 的微波，发射圆天线口径 \(1.5\,\mathrm{km}\)。求中央波束全角宽及地面覆盖圆区的直径（考虑发射天线本身口径）。
\paragraph{习题 23.7}
迎面汽车两灯相距 \(120\,\mathrm{cm}\)，人眼瞳孔直径 \(5.0\,\mathrm{mm}\)，光波长 \(550\,\mathrm{nm}\)。只考虑衍射，求人眼恰可分辨两灯时汽车最大距离。
\paragraph{习题 23.8}
低轨卫星照相机距地面 \(160\,\mathrm{km}\)，要分辨车牌中相距 \(5\,\mathrm{cm}\) 的文字。(1) 求要求的角分辨率；(2) 波长 \(500\,\mathrm{nm}\)，求圆孔镜头最小口径。
\paragraph{习题 23.9}
FAST球面射电望远镜总口径500m，瞬时有效口径300m，工作波长20cm。求瑞利角分辨率。
\paragraph{习题 23.10}
相隔 \(10^4\,\mathrm{km}\) 的两射电望远镜同步记录厘米波信号作长基线干涉，估计角分辨率。可取波长5cm作数值示例，但须保留对波长的依赖。
\paragraph{习题 23.11}
大熊星座 \(\zeta\) 双星两颗恒星角距 \(14''\)，使用 \(\lambda=550\,\mathrm{nm}\) 的光。按瑞利判据求至少所需望远镜口径。
\paragraph{习题 23.12}
双缝间距 \(d=0.10\,\mathrm{mm}\)、各缝宽 \(a=0.02\,\mathrm{mm}\)，\(\lambda=480\,\mathrm{nm}\) 光正入射，其后焦距 \(f=50\,\mathrm{cm}\) 的透镜。求焦平面(1)干涉纹间距，(2)单缝中央衍射亮纹全宽，(3)中央包络内干涉主极大数。
\paragraph{习题 23.13}
光栅宽2.0cm，共刻6000条缝，钠黄光 \(\lambda=589.3\,\mathrm{nm}\) 正入射。求主极大所有允许角位置。
\paragraph{习题 23.14}
光栅每厘米刻6000条，某单色光第1级谱线在 \(20^\circ\)。求该波长和第2级谱线角度。
\paragraph{习题 23.15}
按原教材图23.23的 \(H_\alpha\) 光谱，波长656.3nm的第2级谱线在约 \(41^\circ\)，其第3级光栅谱线缺级。估算栅常 \(d\) 及单缝宽 \(a\)，并保留图示读数误差。
\paragraph{习题 23.16}
氢红双线波长约656.3nm，波长差0.18nm。要求光栅第1级分辨它们，问至少需多少刻线？
\paragraph{习题 23.17}
192个相同天线沿东西方向排成等间距一维阵列，相邻距6m；频率232MHz，馈线相位差为零。求正天顶方向强度极值、向东第一极小角度和下一\emph{非零级主极大}角度。应区别旁瓣的局部极大。
\paragraph{习题 23.18}
晶面间距0.275nm，X射线以掠射角 \(45^\circ\) 照射晶体，光谱涵盖0.095--0.130nm。求可能产生布拉格加强的波长。
\paragraph{习题 23.19}
戴维孙--革末镍晶体实验：电子束垂直入射，沿教材图23.34所定义的 \(\varphi=50^\circ\) 方向出现衍射最大强度，晶面上相邻镍原子间距0.215nm。按表面原子散射的光程差求电子德布罗意波长。
